\documentclass[../../main.tex]{subfiles}
\begin{document}

\begin{center}
  {\fontsize{48pt}{58pt}\selectfont \textbf{2001}\normalsize\textbf{年度}}
\end{center}

\begin{table}[htb]
  \centering
  \caption{出題テーマと難易度}
  \begin{tabular}{|c|p{9cm}|c|}
    \hline
    \textbf{問題番号} & \textbf{テーマ} & \textbf{難易度} \\
    \hline
    第1問 & 回転した接線と三次曲線の交点 & \\
    \hline
    第2問 & 虚軸上に解をもつ五次方程式 & \\
    \hline
    第3問 & 複素数列の周期 & \\
    \hline
    第4問 & 正八面体の頂点と内積の符号 & \\
    \hline
    第5問 & 一の冪根の積が一になる組の個数 & \\
    \hline
    第6問 & 指数関数と振動関数の定積分極限 & \\
    \hline
  \end{tabular}
\end{table}

\begin{table}[htb]
  \centering
  \caption{解答}
  \renewcommand{\arraystretch}{1.4}
  \begin{tabular}{|c|c|p{9cm}|}
    \hline
    \textbf{問題番号} & \textbf{小問} & \textbf{答え} \\
    \hline
    第1問 & -- & $P=(t,t^3)$で$-1/\sqrt{3}<t<1/\sqrt{3}$ \\
    \hline
    第2問 & -- & $a=0,\ \dfrac{1}{\sqrt{2}}-\dfrac{1}{2}$ \\
    \hline
    第3問 & -- & $k$は$8$の正の倍数 \\
    \hline
    第4問 & -- & 条件を満たす頂点$P_k$が存在（証明問題） \\
    \hline
    第5問 & (1) & $a_3=(p-1)(p-2)$ \\
    \cline{2-3}
    & (2) & $a_{n+2}=(p-2)a_{n+1}+(p-1)a_n$ \\
    \cline{2-3}
    & (3) & $a_n=\dfrac{(p-1)^n+(-1)^n(p-1)}{p}$ \\
    \hline
    第6問 & -- & $\dfrac{2}{\pi}(1-e^{-\pi})$ \\
    \hline
  \end{tabular}
\end{table}

\paragraph{試験形態}
6問．試験時間は要確認．

\end{document}
